\documentclass[12pt, letterpaper]{article}

\usepackage[utf8]{inputenc}
\usepackage[T1]{fontenc}
\usepackage[spanish]{babel}
\usepackage{geometry}
\usepackage{amsmath, amssymb, amsthm}
\usepackage{booktabs}
\usepackage{array}
\usepackage{microtype}
\usepackage{setspace}
\usepackage{parskip}
\usepackage{titlesec}
\usepackage{fancyhdr}
\usepackage{enumitem}
\usepackage{bm}
\usepackage{mathtools}
\usepackage{tcolorbox}
\usepackage{hyperref}

\tcbuselibrary{skins, breakable}

\geometry{top=2.8cm, bottom=3.0cm, left=3.2cm, right=3.2cm}
\setstretch{1.20}

\pagestyle{fancy}
\fancyhf{}
\fancyhead[C]{\small\textsc{Econometria · Gujarati --- Ejercicio 11.2}}
\fancyfoot[C]{\small\thepage}
\renewcommand{\headrulewidth}{0.4pt}
\renewcommand{\footrulewidth}{0pt}

\titleformat{\section}
  {\large\bfseries\scshape}{\thesection.}{0.6em}{}
  [\vspace{-4pt}\rule{\linewidth}{0.4pt}]
\titleformat{\subsection}{\normalsize\bfseries}{\thesubsection.}{0.5em}{}
\titleformat{\subsubsection}{\normalsize\itshape}{}{0em}{}

\newtheorem{prop}{Proposicion}
\theoremstyle{remark}
\newtheorem{obs}{Observacion}

\newtcolorbox{clave}{
  colback=white, colframe=black,
  boxrule=0.5pt, arc=3pt,
  left=8pt, right=8pt, top=6pt, bottom=6pt,
  breakable
}

\newcommand{\E}{\mathbb{E}}
\newcommand{\Var}{\operatorname{Var}}
\newcommand{\Cov}{\operatorname{Cov}}

\begin{document}

\begin{center}
  {\LARGE\bfseries Regresion de Salarios sobre Empleados:}\\[6pt]
  {\LARGE\bfseries Heteroscedasticidad, Transformacion WLS}\\[6pt]
  {\LARGE\bfseries y Comparabilidad de $R^2$}\\[12pt]
  {\large\itshape Derivacion completa de la transformacion, supuesto de varianza}\\
  {\large\itshape y limitaciones de la comparacion de coeficientes de determinacion}\\[14pt]
  \rule{0.55\linewidth}{0.6pt}\\[8pt]
  {\normalsize Ejercicio~11.2 --- \textit{Econometria}, Gujarati \& Porter, 5.a ed. por Emanuel Quintana Silva}\\[4pt]
  {\small\textsc{Con algebra de MCP, demostracion de homoscedasticidad del error transformado
  y geometria del $R^2$}}
\end{center}

\vspace{10pt}

\noindent\textbf{Datos del problema.}\quad
Se dispone de una muestra de $n = 30$ empresas con los siguientes resultados:

\begin{align}
  \hat W_i &= 7.5 + 0.009\, N_i
  \tag{1}\label{eq:reg1}\\
  &\quad t: \text{n.a.} \quad (16.10) \qquad R^2 = 0.90 \notag\\[8pt]
  \widehat{W_i/N_i} &= 0.008 + 7.8\,(1/N_i)
  \tag{2}\label{eq:reg2}\\
  &\quad t: (14.43) \quad (76.58) \qquad R^2 = 0.99 \notag
\end{align}

donde $W_i$ es el salario promedio en dolares y $N_i$ es el numero de empleados
de la empresa $i$.

\section{Inciso (a): interpretacion de las dos regresiones}

\subsection{Modelo (1): regresion lineal en niveles}

El modelo subyacente de~( \ref{eq:reg1} ) es:
\begin{equation}
  W_i = \beta_1 + \beta_2 N_i + u_i
  \label{eq:modelo1}
\end{equation}

\begin{itemize}[leftmargin=1.6em, itemsep=5pt]
  \item \textbf{Intercepto $\hat\beta_1 = 7.5$}: Es el salario promedio estimado cuando
    $N_i = 0$. En sentido estricto no tiene interpretacion economica directa
    (una empresa no puede tener cero empleados), pero captura el efecto promedio de
    variables omitidas del modelo.

  \item \textbf{Pendiente $\hat\beta_2 = 0.009$}: Por cada empleado adicional, el
    salario promedio de la empresa aumenta en \$0.009. Equivalentemente, por cada
    1\,000 empleados adicionales el salario medio sube \$9. El estadistico $t$
    asociado es $16.10$, lo que indica que el coeficiente es altamente significativo.

  \item \textbf{$R^2 = 0.90$}: El numero de empleados explica el $90\%$ de la
    variacion del salario promedio entre empresas.
\end{itemize}

\subsection{Modelo (2): regresion en forma transformada}

El modelo subyacente de~(\ref{eq:reg2}) es:
\begin{equation}
  \frac{W_i}{N_i} = \alpha_1 + \alpha_2 \frac{1}{N_i} + v_i
  \label{eq:modelo2}
\end{equation}

\begin{itemize}[leftmargin=1.6em, itemsep=5pt]
  \item \textbf{Variable dependiente $W_i/N_i$}: es el salario total promedio
    \emph{por empleado} de la empresa $i$. Mide la remuneracion unitaria.

  \item \textbf{Regresor $1/N_i$}: es el reciproco del tamano de la empresa.
    Valores grandes de $1/N_i$ corresponden a empresas pequenas y vice-versa.

  \item \textbf{Intercepto $\hat\alpha_1 = 0.008$}: cuando $1/N_i \to 0$ (empresa
    muy grande, $N_i \to \infty$), el salario por empleado converge a \$0.008.
    Es el salario por empleado de largo plazo a medida que la empresa crece.

  \item \textbf{Pendiente $\hat\alpha_2 = 7.8$}: empresas mas pequenas (mayor
    $1/N_i$) tienen mayor salario por empleado. Cada unidad adicional de $1/N_i$
    agrega \$7.8 al salario unitario. Ambos estadisticos $t$ (14.43 y 76.58)
    son altamente significativos.

  \item \textbf{$R^2 = 0.99$}: el modelo (2) explica el 99\% de la variacion de
    $W/N$, un ajuste casi perfecto.
\end{itemize}

\begin{clave}
  El modelo (1) analiza como varia el salario total entre empresas de distinto
  tamano. El modelo (2) analiza como varia el salario por empleado, lo que
  permite comparar empresas en terminos de remuneracion unitaria.
\end{clave}

\section{Inciso (b): supuesto del autor y heteroscedasticidad}

\subsection{La transformacion de (1) a (2)}

El paso de~(\ref{eq:modelo1}) a~(\ref{eq:modelo2}) se obtiene dividiendo cada
termino de~(\ref{eq:modelo1}) entre $N_i$:
\begin{equation}
  \frac{W_i}{N_i} = \frac{\beta_1}{N_i} + \beta_2 + \frac{u_i}{N_i}
  \label{eq:division}
\end{equation}
Reordenando:
\begin{equation}
  \frac{W_i}{N_i} = \beta_2 + \beta_1 \cdot \frac{1}{N_i} + v_i,
  \qquad v_i = \frac{u_i}{N_i}
  \label{eq:transformado}
\end{equation}

Nótese que los roles de $\beta_1$ y $\beta_2$ se intercambian: el intercepto del
modelo original ($\beta_1$) se convierte en la pendiente del modelo transformado
(coeficiente de $1/N_i$), y la pendiente original ($\beta_2$) se convierte en el
intercepto del modelo transformado.

\subsection{El supuesto de heteroscedasticidad implicito}

La transformacion~(\ref{eq:division}) es una tecnica de \textbf{Minimos Cuadrados
Ponderados (MCP)} que el autor utiliza para corregir heteroscedasticidad. El
supuesto implicito sobre la estructura de la varianza del error es:

\begin{equation}
  \Var(u_i) = \sigma^2 N_i^2
  \label{eq:supuesto_heterosc}
\end{equation}

Es decir, la varianza del error en la ecuacion original es proporcional al
cuadrado del numero de empleados. Esto tiene sentido economico: empresas mas
grandes tienen mayor diversidad interna en sus estructuras salariales, por lo que
la variacion alrededor de la relacion media es mayor para empresas con muchos
empleados.

\subsection{Verificacion: el error transformado es homoscedastico}

Bajo el supuesto~(\ref{eq:supuesto_heterosc}), la varianza del error transformado
$v_i = u_i/N_i$ es:
\begin{equation}
  \Var(v_i) = \Var\left(\frac{u_i}{N_i}\right)
            = \frac{1}{N_i^2}\Var(u_i)
            = \frac{1}{N_i^2} \cdot \sigma^2 N_i^2
            = \sigma^2
  \label{eq:homosc_vi}
\end{equation}

La varianza de $v_i$ es constante e igual a $\sigma^2$: el error transformado
es \textbf{homoscedastico}. Por ello, aplicar MCO a~(\ref{eq:transformado}) es
equivalente a aplicar MCP a~(\ref{eq:modelo1}) con pesos $w_i = 1/N_i^2$.
Los estimadores resultantes son MELI (eficientes) bajo el supuesto~(\ref{eq:supuesto_heterosc}).

\subsection{Evidencia de que al autor le preocupaba la heteroscedasticidad}

Se sabe que al autor le preocupaba la heteroscedasticidad por dos razones:

\begin{enumerate}[leftmargin=1.8em, itemsep=4pt]
  \item La transformacion division por $N_i$ es una tecnica remedial clasica
    cuando se sospecha que $\Var(u_i) \propto N_i^2$. No habria razon para
    aplicarla si no se sospechara varianza no constante.

  \item El modelo (2) tiene $R^2 = 0.99$ frente al $0.90$ del modelo (1),
    y los estadisticos $t$ del modelo (2) son mucho mas grandes (76.58 vs 16.10),
    lo cual es consistente con la ganancia de eficiencia que se obtiene al
    corregir la heteroscedasticidad.
\end{enumerate}

\begin{clave}
  El supuesto de heteroscedasticidad es $\Var(u_i) = \sigma^2 N_i^2$. Al dividir
  por $N_i$, el error transformado $v_i = u_i/N_i$ tiene varianza constante
  $\sigma^2$, restaurando las propiedades MELI de los estimadores MCO aplicados
  al modelo transformado.
\end{clave}

\section{Inciso (c): relacion entre pendientes e interceptos de los dos modelos}

\subsection{Relacion algebraica exacta}

Comparando el modelo original~(\ref{eq:modelo1}) con el transformado~(\ref{eq:transformado}):

\begin{center}
\begin{tabular}{lcc}
\toprule
\textbf{Parametro} & \textbf{Modelo (1)} & \textbf{Modelo (2)} \\
\midrule
Intercepto & $\beta_1$ & $\alpha_1 = \beta_2$ \\
Pendiente  & $\beta_2$ & $\alpha_2 = \beta_1$ \\
\bottomrule
\end{tabular}
\end{center}

\vspace{6pt}

Es decir, \textbf{el intercepto del modelo (1) se convierte en la pendiente del
modelo (2), y la pendiente del modelo (1) se convierte en el intercepto del
modelo (2)}. La relacion exacta es:

\begin{align}
  \hat\alpha_2^{(2)} &= \hat\beta_1^{(1)} \quad \text{(intercepto de (1) = pendiente de (2))}
  \label{eq:rel1}\\
  \hat\alpha_1^{(2)} &= \hat\beta_2^{(1)} \quad \text{(pendiente de (1) = intercepto de (2))}
  \label{eq:rel2}
\end{align}

\subsection{Valores numericos y discrepancia}

\begin{center}
\small
\begin{tabular}{lcc}
\toprule
& \textbf{Modelo (1) [MCO]} & \textbf{Modelo (2) [MCP]} \\
\midrule
$\hat\beta_1$ (intercepto) & $7.5$ & --- \\
$\hat\beta_2$ (pendiente)  & $0.009$ & --- \\
$\hat\alpha_2$ (pendiente de $1/N$) & --- & $7.8$ \\
$\hat\alpha_1$ (intercepto) & --- & $0.008$ \\
\midrule
\multicolumn{3}{l}{Comparacion: $7.5 \approx 7.8$ y $0.009 \approx 0.008$} \\
\bottomrule
\end{tabular}
\end{center}

\vspace{6pt}

Los valores son aproximadamente iguales pero no identicos ($7.5 \neq 7.8$,
$0.009 \neq 0.008$). La discrepancia ocurre porque:

\begin{enumerate}[leftmargin=1.8em, itemsep=4pt]
  \item El modelo (1) se estima por MCO ordinario, que asume homoscedasticidad
    y asigna igual peso a todas las observaciones.
  \item El modelo (2) se estima por MCO sobre las variables transformadas, lo
    que es equivalente a MCP con pesos $w_i = 1/N_i^2$. Las empresas grandes
    (mayor $N_i$) reciben menor peso, lo que cambia los estimadores.
\end{enumerate}

Si el supuesto~(\ref{eq:supuesto_heterosc}) es correcto, los estimadores del
modelo (2) son mas eficientes que los del modelo (1), por lo que las estimaciones
de (2) son mas precisas aunque difieran numericamente de las de (1).

\begin{clave}
  Algebraicamente: pendiente de (2) $= \hat\alpha_2 \approx \hat\beta_1$ y
  intercepto de (2) $= \hat\alpha_1 \approx \hat\beta_2$. La pequena discrepancia
  numerica refleja la diferente ponderacion de las observaciones en MCO vs MCP.
\end{clave}

\section{Inciso (d): ¿se pueden comparar los $R^2$ de los dos modelos?}

\subsection{Respuesta}

\textbf{No}, los $R^2$ de los dos modelos \textbf{no son comparables directamente}.

\subsection{Razon tecnica}

El $R^2$ se define como la proporcion de la variacion total de la variable
dependiente que es explicada por el modelo:
\begin{equation}
  R^2 = 1 - \frac{\text{SCR}}{\text{SCT}}
      = 1 - \frac{\sum_i (Y_i - \hat Y_i)^2}{\sum_i (Y_i - \bar Y)^2}
  \label{eq:R2_def}
\end{equation}

Las variables dependientes de los dos modelos son distintas:
\begin{align}
  \text{Modelo (1):} \quad &Y = W_i \quad \Rightarrow \quad \text{SCT}_1 = \sum(W_i - \bar W)^2\\
  \text{Modelo (2):} \quad &Y = W_i/N_i \quad \Rightarrow \quad \text{SCT}_2 = \sum(W_i/N_i - \overline{W/N})^2
\end{align}

Como SCT$_1 \neq$ SCT$_2$, los valores de $R^2$ de los dos modelos miden
fracciones de totales completamente distintos. Concretamente:

\begin{itemize}[leftmargin=1.6em, itemsep=4pt]
  \item $R^2_1 = 0.90$ dice que el modelo (1) explica el 90\% de la variacion de
    los salarios medios totales $W_i$.
  \item $R^2_2 = 0.99$ dice que el modelo (2) explica el 99\% de la variacion de
    los salarios por empleado $W_i/N_i$.
\end{itemize}

Que $R^2_2 > R^2_1$ no implica que el modelo (2) sea ``mejor'' en algun sentido
absoluto: simplemente explica una proporcion mayor de una variable que ha sido
escalada de forma diferente. Un $R^2 = 0.99$ en la regresion de $W/N$ no puede
compararse con un $R^2 = 0.90$ en la regresion de $W$ del mismo modo que no
tendria sentido comparar el $R^2$ de una regresion en dolares con el $R^2$ de
la misma regresion en centavos.

\subsection{Condicion formal de no comparabilidad}

Una condicion suficiente para que dos $R^2$ sean comparables es que las
variables dependientes sean identicas (o al menos que sean transformaciones
monotónicas que preserven los ordenes de magnitud). Formalmente, si
$Y^{(1)} \neq Y^{(2)}$, entonces:
\begin{equation}
  R^2_1 - R^2_2
    = \frac{\text{SCR}_2}{\text{SCT}_2} - \frac{\text{SCR}_1}{\text{SCT}_1}
  \label{eq:dif_R2}
\end{equation}
no tiene interpretacion directa porque los denominadores son incomensurables.

\begin{clave}
  Los $R^2$ de los dos modelos \textbf{no son comparables} porque las variables
  dependientes son diferentes ($W$ vs $W/N$). El $R^2$ mide la fraccion de la
  variacion de su propia variable dependiente que el modelo explica; como esas
  variaciones totales son distintas, los cocientes no son comparables entre si.
\end{clave}

\section{Sintesis}

\begin{table}[ht]
\centering
\caption{Comparacion entre los modelos (1) y (2)}
\label{tab:comparacion}
\small
\begin{tabular}{p{3.5cm} p{5cm} p{5cm}}
\toprule
\textbf{Aspecto}
  & \textbf{Modelo (1): $W = \beta_1 + \beta_2 N + u$}
  & \textbf{Modelo (2): $W/N = \alpha_1 + \alpha_2/N + v$} \\
\midrule
Variable dependiente
  & Salario total promedio $W$
  & Salario por empleado $W/N$ \\[4pt]
Regresor
  & Numero de empleados $N$
  & Reciproco del tamano $1/N$ \\[4pt]
Metodo de estimacion
  & MCO ordinario
  & MCO sobre datos transformados (= MCP) \\[4pt]
Supuesto de varianza
  & Homoscedasticidad $\Var(u_i) = \sigma^2$
  & Heteroscedasticidad $\Var(u_i) = \sigma^2 N_i^2$ \\[4pt]
Eficiencia estimadores
  & MELI solo bajo homoscedasticidad
  & MELI bajo $\Var(u_i) = \sigma^2 N_i^2$ \\[4pt]
$R^2$
  & $0.90$ (variacion de $W$)
  & $0.99$ (variacion de $W/N$) \\[4pt]
Comparabilidad de $R^2$
  & \multicolumn{2}{c}{No comparables (variables dependientes distintas)} \\
\bottomrule
\end{tabular}
\end{table}

\appendix

\section{Demostracion de que MCP con pesos $1/N_i^2$ es equivalente a MCO sobre las variables transformadas}

En el modelo de MCP se minimiza la suma ponderada:
\begin{equation}
  \min_{\beta_1,\beta_2} \sum_i w_i (W_i - \beta_1 - \beta_2 N_i)^2,
  \qquad w_i = \frac{1}{N_i^2}
  \label{eq:MCP}
\end{equation}

Multiplicando y dividiendo por $N_i$ dentro del cuadrado:
\begin{equation}
  \sum_i \frac{1}{N_i^2}(W_i - \beta_1 - \beta_2 N_i)^2
  = \sum_i \left(\frac{W_i}{N_i} - \frac{\beta_1}{N_i} - \beta_2\right)^2
  \label{eq:MCP_equiv}
\end{equation}

Definiendo $Y_i^* = W_i/N_i$, $X_i^* = 1/N_i$ y $\alpha_1 = \beta_2$, $\alpha_2 = \beta_1$:
\begin{equation}
  \sum_i (Y_i^* - \alpha_1 - \alpha_2 X_i^*)^2
  \label{eq:MCO_transformado}
\end{equation}

Minimizar~(\ref{eq:MCO_transformado}) por MCO ordinario es exactamente
equivalente a minimizar~(\ref{eq:MCP}). Los estimadores MCO de $\alpha_1$ y
$\alpha_2$ son los mismos que los estimadores MCP de $\beta_2$ y $\beta_1$,
respectivamente. Esto confirma la relacion algebrica del inciso~(c).

\section{Generalidad del supuesto de varianza proporcional}

El supuesto $\Var(u_i) = \sigma^2 N_i^2$ es un caso particular del supuesto
general $\Var(u_i) = \sigma^2 X_i^{2p}$. Para distintos valores de $p$:

\begin{center}
\small
\begin{tabular}{cll}
\toprule
$p$ & $\Var(u_i)$ & Transformacion para homoscedasticidad \\
\midrule
$0$ & $\sigma^2$ (homoscedastico) & Sin transformacion \\
$1/2$ & $\sigma^2 X_i$ & Dividir por $\sqrt{X_i}$ \\
$1$ & $\sigma^2 X_i^2$ & Dividir por $X_i$ (caso del ejercicio) \\
$3/2$ & $\sigma^2 X_i^3$ & Dividir por $X_i^{3/2}$ \\
\bottomrule
\end{tabular}
\end{center}

\noindent En el ejercicio, el autor elige $p = 1$ (dividir por $N_i$), lo que
implica el supuesto $\Var(u_i) = \sigma^2 N_i^2$. Si el supuesto real fuera
$\Var(u_i) = \sigma^2 N_i$, la transformacion correcta seria dividir por
$\sqrt{N_i}$, produciendo un modelo diferente.

\vspace{14pt}
\noindent\rule{\linewidth}{0.4pt}

\noindent{\small\textit{Nota.} La no comparabilidad de los $R^2$ es un punto
tecnico importante en econometria aplicada. Siempre que se compare el ajuste
de modelos con distintas variables dependientes (niveles vs logaritmos, totales
vs per capita, etc.), el $R^2$ convencional no es un criterio valido de
comparacion. En esos casos se usan criterios alternativos como el AIC, el BIC,
o el $R^2$ de Theil ajustado para transformaciones.}

\end{document}
